\subsection{Herausforderungen und Grenzen von KI in ERP-Systemen}
Trotz der vielversprechenden Anwendungsfälle und Potenziale von \ac{KI} in \ac{ERP}-Systemen besteht auch eine Reihe von Herausforderungen und Grenzen. 

Zunächst ist die Einbettung von \ac{KI}-Technologien in bestehende \ac{ERP}-Systeme technisch komplex, da \ac{ERP}-Systeme aus Millionen von Codezeilen bestehen und verschiedenste branchenspezifische sowie regionale Anforderungen erfüllen müssen. \cite[\pagef 73342]{sarferaz_implementing_2025}

Darüber hinaus stellt die Erstellung generativer Ergebnisse durch \ac{KI}-Modelle eine wesentliche Schwäche dar, da diese gelegentlich ungenaue oder unpassende Informationen erzeugen können. Dies macht manuelle Eingriffe und Korrekturen erforderlich, wodurch das Vertrauen in das System beeinträchtigt werden kann. \cite[\pagef 73344f]{sarferaz_implementing_2025}

Auch bei der Nutzung von \ac{KI} in Form von \ac{KI}-Agenten bestehen Herausforderungen. Ein zentrales Problem stellt die Qualität der zugrunde liegenden Daten dar. KI-Agenten sind in hohem Maße auf konsistente, vollständige und aktuelle Daten angewiesen, um verlässliche Entscheidungen treffen zu können. In bestehenden \ac{ERP}-Landschaften sind Daten jedoch häufig historisch gewachsen, fragmentiert oder inkonsistent, was die Leistungsfähigkeit \ac{KI}-gestützter Prozesse erheblich einschränken kann. Fehlerhafte oder verzerrte Daten können dabei direkt zu unzutreffenden Handlungsempfehlungen oder fehlerhaften automatisierten Aktionen führen. \cite[\pagef 1035]{gujar_role_2025}

Darüber hinaus ist die Integration von \ac{KI}-Agenten in bestehende \ac{ERP}-Systeme mit erheblichen technischen Herausforderungen verbunden. Insbesondere die enge Kopplung geschäftskritischer Prozesse sowie der Zugriff auf sensible Unternehmens- und Kundendaten erfordern hohe Anforderungen an Datensicherheit und Datenschutz. Der Einsatz von \ac{KI}-Agenten muss daher mit bestehenden Sicherheitskonzepten und regulatorischen Vorgaben in Einklang gebracht werden, um Risiken für Datenschutzverletzungen oder unbefugte Systemzugriffe zu vermeiden. \cite[\pagef 1035f]{gujar_role_2025}

Zusätzlich führt die Einführung agentenbasierter Architekturen zu einer erhöhten Systemkomplexität innerhalb bestehender \ac{ERP}-Systeme. Zusätzliche Komponenten erhöhen den administrativen und operativen Aufwand und stellen zusätzliche Anforderungen an Betrieb und Wartung der Systeme. \cite[\pagef 178956f]{sarferaz_implementing_2025-1}

Ein weiterer limitierender Faktor betrifft die Performance von \ac{KI}-Agenten im produktiven \ac{ERP}-Betrieb. Experimentelle Ergebnisse zeigen, dass der Einsatz von \ac{KI}-Agenten mit zusätzlichem Speicherbedarf sowie einem messbaren CPU-Overhead verbunden ist. \cite[\pagef 178956f]{sarferaz_implementing_2025-1}

